\chapter{CONCLUSION}\label{ch:conclusion}
\begingroup\emergencystretch=3em
\noindent VeerDrishti has progressed from a proposed combat-vision framework to an Android prototype that
connects camera capture, on-device inference, and visual output. The demonstration identifies
CameraX, YOLO26 Nano task variants, Flutter YOLOView, and Logcat as the
main implementation components. The presented screenshots demonstrate object detection,
segmentation, and multi-person pose overlays in portrait and landscape views, with configurable
thresholds and runtime indicators. This establishes the application's capture-to-display
foundation and provides concrete evidence for the Implementation and Result chapters.

The observed runtime readings and visual exhibits do not establish battlefield accuracy,
sustained performance, or a reduction in casualties. Those outcomes remain project objectives
that require combat-relevant data and controlled evaluation. The next development steps are to
document the deployed configuration, evaluate accuracy and latency on repeatable test cases,
refine the interface, and validate the selected model for the intended object classes. Thermal
fusion, open-vocabulary recognition, and persistent event recording remain possible extensions
rather than demonstrated features. By grounding further development in measured evidence,
VeerDrishti can advance toward a more reliable, portable visual decision-support system.
\label{end:conclusion}
\par\endgroup
